%% Alternative back page, written by Bernd Paysan
%% Translated into Spanish (es_ES) by Francisco Escobedo
%% you'll find the original back page at version 1.3 in the CVS repository

\noindent
\textbf{Pensando en \Forth{}} es un libro sobre la filosof�a de
soluci�n de problemas y estilo de programaci�n, aplicados al lenguaje
�nico de programaci�n `\Forth{}'. Publicado originalmente en 1984,
puede estar entre los libros cl�sicos atemporales de la Inform�tica,
como \emph{The Mythical Man-Month}, de \person{Fred Brooks}, y
\emph{The Art of Computer Programming}, de \person{Donald Knuth}.

Muchos principios de ingenier�a de la programaci�n descritos aqu� han
sido redescubiertos en la Programaci�n eXtrema, incluyendo el
(re)descomponer, la modularidad y el dise�o ascendente y paulatino.
Aqu� encontrar�s todo eso y m�s, tal como el valor del an�lisis y el
dise�o, descrito en el estilo sensato y humor�stico de \person{Leo
  Brodie}, con ilustraciones, listados, aplicaciones pr�cticas del
mundo real, dibujos explicativos y entrevistas con el inventor de
\Forth{}, \person{Charles H. Moore}, as� como otros pensadores en
\Forth{}.

Si programa en \Forth{}, este es un libro imprescindible. Si no, los
conceptos fundamentales son universales: \textbf{Pensando en \Forth{}}
est� dirigido a cualquiera que est� interesado en escribir programas
para resolver problemas. Los conceptos van m�s all� de \Forth{}, pero
la \emph{simple belleza} de \Forth{} revela tales conceptos con total
claridad.

As� que abra el libro y l�alo todo sobre la filosof�a de \Forth{}, el
an�lisis, la descomposici�n, la resoluci�n de problemas, el estilo y
los convenios, la factorizaci�n, la manipulaci�n de datos y la
minimizaci�n de estructuras de control. Pero tenga cuidado: tal vez no
pueda dejarlo una vez comenzado.

\medskip
Este libro ha sido digitalizado, pasado por OCR, maquetado en
\LaTeX{}, y llevado de nuevo a la imprenta (y el monitor) por un
esfuerzo de colaboraci�n bajo una licencia Creative Commons.

\medskip\blackline{0pt}\medskip
%%if you want Leo's picture here, uncommend this
\ifleo
\makeatletter
\setbox\@tempboxa\hbox{\includegraphics{leobrodie}}
\rightmargin\wd\@tempboxa
\@tempdima\ht\@tempboxa
\advance\@tempdima-\ht\strutbox
\advance\rightmargin1em
\advance\textwidth-\rightmargin
\makeatother
\noindent\begin{tabular}{@{}p{\textwidth}@{\hspace{1em}}l@{}}
\parindent2.5em
\fi
\textbf{\person{Leo Brodie}} ha pensado profesionalmente sobre el
desarrollo de programas desde 1980, cuando entr� en \Forth{}, Inc.\@
como escritor t�cnico. Tambi�n ha trabajado como asesor, autor,
desarrollador y director de programas para empresas como IBM, NCR,
Digalog Corp., Microsoft, y Real Networks. Tambi�n es autor de
\emph{Starting \Forth{}} (Prentice-Hall, 1981).
\ifleo
& \makeatletter
\smash{\lower\@tempdima\hbox{\includegraphics{leobrodie}}}
\makeatother \\
\end{tabular}
\fi